\documentclass[aspectratio=169]{beamer}

%\usepackage[dark, english]{SU-Ray} 可选黑色模板/option for dark mode
%\usepackage[swedish]{SU-Ray} 瑞典语版logo致谢词/Logo and thanks in Swedish
\usepackage[swedish]{SU-Ray} 

%% ==========================================
%% 首页信息填空区 / Title Page Information Area
%% ==========================================
\title{An Extremely Long and Comprehensive Guide to Modern LaTeX Presentation Design Principles} 
\subtitle{Exploring Advanced Beamer Features, Custom Color Palettes, and Elegant Typography for Academic Defenses}

\author{Ray SU}                             
\degreeproject{Master}{Data Science, Statistics and Decision Analysis}       
\supervisor{Supervisor Name}                    
\institute{Department of Computer and Systems Sciences} 
\date{June 2026}                                  
%% ==========================================

\begin{document}

%% --- 首页 /  Title Page ---
\SUtitlepage 

%% --- 目录页 /  Table of Contents ---
%% \Sutocpage[Content] []里写任何你想替代outline的东西
%% \Sutocpage[Content] []fill in anything you want to replace Outline
\SUtocpage

%% ==========================================
%% Section 1: 逐行展示与基础排版
%% Section 1: Step-by-Step Animations & Basic Layout
%% ==========================================
\section{Step-by-Step Animations}

\begin{frame}{Dynamic Lists}
    % \pause 会让内容在放映时分步出现
    % \pause makes content appear step-by-step during presentation
    Welcome to the SU-Ray Beamer Theme. Let's look at the basic features:
    \vspace{0.3cm}
    \begin{itemize}
        \item This is the first bullet point. \pause
        \item Notice how the active items use the primary \alert{SUblue} color for bullets. \pause
        \item You can emphasize text using the \textbackslash alert command, which renders in \alert{SUfire} orange. \pause
    \end{itemize}
    
    \vspace{0.3cm}
    \begin{alertblock}{Important Note}
        Smooth transitions keep the audience focused on your current talking point.
    \end{alertblock}
\end{frame}

%% ==========================================
%% Section 2: 颜色与专属区块
%% Section 2: Colors and Blocks
%% ==========================================
\section{Colors and Blocks}

\begin{frame}{Information Blocks}
    % 标准区块 (官方大学蓝) / Standard Block (Official SU Blue)
    \begin{block}{Standard Block}
        This block is used for general information, definitions, or standard theorems. 
    \end{block}
    
    % 示例区块 (学术橄榄绿) / Example Block (Academic Olive Green)
    \begin{exampleblock}{Example Block}
        Perfect for case studies, examples, or practical applications.
    \end{exampleblock}
    
    % 警告区块 (火橙色) / Alert Block (Fire Orange)
    \begin{alertblock}{Alert Block}
        Highlights critical issues, warnings, or extremely important results.
    \end{alertblock}
\end{frame}

%% ==========================================
%% Section 3: 图文双栏排版
%% Section 3: Image & Column Layout
%% ==========================================
\section{Image \& Column Layout}

\begin{frame}{Two-Column Layout}
    \begin{columns}[T] % T 代表顶部对齐 / T means top-aligned
        
        % 左侧栏：占屏幕宽度的 50% / Left column: 50% width
        \begin{column}{0.5\textwidth}
            \textbf{Left Column Text} \\
            \vspace{0.5cm}
            You can easily split the frame into two columns for better space utilization on a 16:9 screen. 
            \vspace{0.5cm}
            \begin{itemize}
                \item Clean visual separation
                \item Ideal for comparing data
                \item Text next to images
            \end{itemize}
        \end{column}
        
        % 右侧栏：占屏幕宽度的 45% / Right column: 45% width (leaves 5% gap for breathing room)
        \begin{column}{0.45\textwidth}
            \textbf{Right Column Image} \\
            \vspace{0.5cm}
            \begin{center}
                \includegraphics[width=0.9\textwidth]{Stockholm.jpg}
                \footnotesize \textcolor{SUgrey}{\textit{Figure 1: Placeholder for an actual chart.}}
            \end{center}
        \end{column}
        
    \end{columns}
\end{frame}

%% ==========================================
%% Section 4: 演示多级抬头与长公式分页
%% Section 4: Subsections and Formula Pagination
%% ==========================================
\section{Mathematical Derivation}

%% 声明 Subsection 后，接下来的页面抬头会自动变成 "4.1 xxx"
%% After declaring a Subsection, the following frame titles will be "4.1 xxx"
\subsection{The Normal Distribution}

% 加入 [allowframebreaks] 后，超出页面的内容会自动切分到下一页，并加上罗马数字 (I, II)
% With [allowframebreaks], overflowing content automatically splits to the next page with Roman numerals (I, II)
\begin{frame}[allowframebreaks]{Gaussian Probability Density}
    
    In probability theory, a normal distribution is a type of continuous probability distribution for a real-valued random variable.
    
    \vspace{0.5cm}
    \begin{block}{Probability Density Function}
        The general form of its probability density function is:
        $$ f(x) = \frac{1}{\sigma \sqrt{2\pi}} e^{-\frac{1}{2}\left(\frac{x-\mu}{\sigma}\right)^2} $$
    \end{block}
    
    \vspace{0.5cm}
    The parameter $\mu$ is the mean or expectation of the distribution, while the parameter $\sigma$ is its standard deviation. The variance of the distribution is $\sigma^2$.
    
    %% \framebreak 强制切分到下一页 / Forces a split to the next page
    \framebreak 
    
    A standard normal distribution has a mean of 0 and a standard deviation of 1. The Gaussian integral is a fundamental property:
    
    \vspace{0.5cm}
    \begin{block}{Gaussian Integral}
        $$ \int_{-\infty}^{\infty} e^{-x^2} dx = \sqrt{\pi} $$
    \end{block}
    
    \vspace{0.5cm}
    Notice how the frame title automatically appended a "II" to indicate this is the continuation!

\end{frame}

%% ==========================================
%% Section 5: 代码展示
%% Section 5: Code Integration
%% ==========================================
\section{Implementation}
\subsection{Basic Syntax}

%% 注意：带有代码的 frame 要加 [fragile] 参数！
\begin{frame}[fragile]{C Language Code Snippet}
    
    To insert source code perfectly into Beamer, always remember to add \\the \textbf{[fragile]} option to the frame environment.
    
    \vspace{0.3cm}
    % 使用我们新定义的 macode 环境！
    % 中括号 [] 里写语言，大括号 {} 里写窗口标题(文件名)
    \begin{macode}[language=C]{hello\_world.c}{Show mac code}
#include <stdio.h>

/* This is a comment: The main function */
int main() {
    // Print a greeting to the console
    printf("Hello, World!\n");
    
    return 0;
}
    \end{macode}
    
\end{frame}

%% ==========================================
%% Section 6: Bonus (隐藏彩蛋)
%% Section 6: Bonus (Easter Egg)
%% ==========================================
\section{Bonus}

\begin{frame}{A Well-Deserved Break}
    \begin{center}
        \includegraphics[height=4.5cm]{semla.png}\par
        \vspace{0.6cm}
        \textcolor{SUblue}{\Large \textit{"Any academic presentation longer than 20 minutes requires a mandatory Swedish Fika."}}
        
    \end{center}
\end{frame}

%% --- 尾页致谢 /  Thank You Page ---
\SUthankyou 

\end{document}